% part: intuitionistic-logic
% chapter: soundness-completeness
% section: soundness-axd

\documentclass[../../../include/open-logic-section]{subfiles}

\begin{document}

\olfileid{int}{sc}{sax}

\olsection{د بديهي \usetoken{P}{derivation} سموالے}

\begin{editorial}
  په بديهي حساب کښې هره کرښه يا مقدمه ده، يا د بديهي بڼې بېلګه،
  يا له دوو پخوانيو کرښو د مودوس پوننس له مخې راځي۔
\end{editorial}

\begin{thm}[سموالے]
  \ollabel{thm:soundness}
  که $\Gamma \Proves !A$ وي، نو $ \Gamma \Entails !A$ وي۔
\end{thm}

\begin{proof}
  ثابتوو چې که $\Gamma \Proves !A$ وي، نو $\Gamma \Entails !A$
  وي۔ ثبوت له~$\Gamma$ څخه د~$!A$ په !!{derivation} کښې د
  !!{formula}s پر شمېر~$n$ استقرا ده۔ ښيو چې که $!A_1$،
  \dots، $!A_n = !A$ له~$\Gamma$ څخه !!a{derivation} وي، نو
  $\Gamma \Entails !A_n$ وي۔ پام وکړئ چې که $!A_1$، \dots،
  $!A_n$ !!a{derivation} وي، نو د هر $k<n$ لپاره $!A_1$،
  \dots، $!A_k$ هم استنتاج دے۔

  د اوږدوالي~$0$ هېڅ !!{derivation}s نشته، نو د $n=0$ لپاره ادعا
  تشه رښتيا ده۔ اوس فرض کوو چې ادعا د اوږدوالي~$<n$ د ټولو
  !!{derivation}s لپاره سمه ده۔ د~$!A_n$ د توجيه له مخې حالتونه
  بېلوو۔
  
  \begin{enumerate}
  \item $!A_n$ بديهي اصل دے۔ ټول بديهي اصول معتبر دي، نو
    د هر~$\Gamma$ لپاره $\Gamma \Entails !A_n$ وي۔
    د عطف لومړۍ دوې بڼې د جوړې غړي راباسي۔ د عطف درېيمه بڼه
    په دوو پرله‌پسې د لاسرسي وړ نړۍو کښې دواړه فرضونه سره يوځاے
    کوي، ځکه لومړے فرض تر دويمې نړۍ پورې ساتل کېږي۔ د فصل لومړۍ
    دوې بڼې د اړوند غړي له مخې فصل رښتينے کوي۔ د فصل درېيمه
    بڼه په وروستۍ نړۍ کښې پر دواړو شرطونو عمل کوي، او د فصل
    له رښتيني غړي څخه پايله اخلي۔ د شرط په لومړۍ بڼه کښې
    لومړے فرض په ټولو ورپسې نړۍو کښې پاتې کېږي۔ په دويمه
    بڼه کښې په هره وروستۍ نړۍ کښې د لومړي فرض په منلو سره
    د دويم فرض او ورپسې پايلې تر منځ شرط لرو۔ د بلې مقدمې
    له مخې دويم فرض هم لرو، او انعکاسيت دواړه شرطونه په
    هماغې نړۍ کښې تطبيقوي۔ د محال لپاره هېڅ نړۍ مقدمه نۀ
    رښتيا کوي، نو د شرطي بڼې عمومي حکم پوره کېږي۔
  \item $!A_n \in \Gamma$ وي۔ نو د هر $\mModel{M}$ او $w$
    لپاره، که $\mSat{M}{\Gamma}[w]$ وي، څرګنده ده چې
    $\mSat{M}{!A_n}[w]$ وي؛ يعنې $\Gamma \Entails !A$ وي۔
  \item $!A_n$ له $!A_i$ او $!A_j \ident !A_i \lif !A_n$ څخه د
    \MP{} له مخې راځي۔ $!A_1$، \dots، $!A_i$ او $!A_1$،
    \dots، $!A_j$ له~$\Gamma$ څخه !!{derivation}s دي، نو د
    استقرايي فرض له مخې $\Gamma \Entails !A_i$ او
    $\Gamma \Entails !A_i \lif !A_n$ وي۔
    
    فرض کړئ $\mSat{M}{\Gamma}[w]$ وي۔ ځکه
    $\mSat{M}{\Gamma}[w]$ او $\Gamma \Entails !A_i \lif !A_n$
    دي، نو $\mSat{M}{!A_i \lif !A_n}[w]$ وي۔ د تعريف له مخې
    د هر داسې $w'$ لپاره چې $Rww'$ وي، که
    $\mSat{M}{!A_i}[w']$ وي، نو $\mSat{M}{!A_n}[w']$ وي۔
    څرنګه چې $R$ انعکاسي ده، $w$ هم په هغو $w'$ کښې دے چې
    $Rww'$ لري؛ يعنې که $\mSat{M}{!A_i}[w]$ وي، نو
    $\mSat{M}{!A_n}[w]$ وي۔ ځکه چې $\Gamma \Entails !A_i$،
    $\mSat{M}{!A_i}[w]$ وي۔ نو $\mSat{M}{!A_n}[w]$ شو،
    لکه چې غوښتل مو۔
  \end{enumerate}
\end{proof}

\end{document}
